\appendix

\section{Netlist do circuito}
\label{ap:netlist}

As netlists completas das variantes são geradas automaticamente a
partir dos pesos treinados e residem em
\texttt{experiments/analog-mlp/}: \texttt{analog\_mlp.cir}
(amplificador ideal), \texttt{analog\_mlp\_741.cir} (LM741),
\texttt{analog\_mlp\_moderno.cir} (LT1810) e as versões de janela
reduzida \texttt{analog\_mlp\_moderno\_fast.cir} e
\texttt{analog\_mlp\_moderno\_comp.cir}
(Seção~\ref{sec:moderno-rede}), com cerca de 6600 linhas cada.
Abaixo, um trecho representativo,
exatamente como simulado: trilhos, um pixel de entrada com o seu
trilho negado, um neurônio oculto completo e os comparadores.

{\small
\begin{verbatim}
* Analog MLP - 5-stage op-amp MNIST classifier
.include claw_opamps.lib
VCC vcc 0 5
VEE vee 0 -5
VEER veer 0 -0.2
VBP bp 0 1
VBN bn 0 -1
VTH vth 0 2.65

* entrada: um pixel (fonte PWL) e o seu trilho negado
VX10 x10 0 PWL(0 0.0479 250u 0.0479 260u 0 500u 0 ...)
BNX10 nx10 0 V=-V(x10)

* ---- stage 4: 24 -> 16 (hidden, single-supply ReLU)
XS4_0 0 s4_0 vcc veer a4_0 CLAW_IDEAL_OPAMP
RF4_0 a4_0 s4_0 100k
R4_0_0 na3_0 s4_0 619k
R4_0_1 a3_1 s4_0 649k
   ... (24 ramos; peso positivo entra pelo trilho negado na3_i)
RB4_0 bp s4_0 866k
XI4_0 0 si4_0 vcc vee na4_0 CLAW_IDEAL_OPAMP
RIA4_0 a4_0 si4_0 100k
RIF4_0 na4_0 si4_0 100k

* ---- comparators: one op-amp per digit, fires above VTH
XC0 z0 vth vcc veer out0 CLAW_IDEAL_OPAMP
RL0 out0 0 100k
   ...
.meas TRAN z_v0_d0 FIND V(z0) AT=200u
.meas TRAN out_v0_d0 FIND V(out0) AT=200u
   ... (400 medidas)
.tran 0 5000u 0 2u
.end
\end{verbatim}
}

Nas variantes com amplificador comercial mudam apenas a biblioteca
(\texttt{models/lm741.lib} para o LM741, obtida à parte por não ser
redistribuível, e \texttt{LTC.lib}, embarcada no LTspice, para o
LT1810) e os trilhos de alimentação, calibrados com
\texttt{bench\_741\_sat.cir} e \texttt{bench\_moderno\_sat.cir} para
posicionar a saturação dos neurônios ocultos em \SI{0}{\volt} e
\SI{4,8}{\volt}. O macromodelo ideal da primeira variante:

{\small
\begin{verbatim}
.subckt CLAW_IDEAL_OPAMP INP INN VCC VEE OUT params: Aol=100k Rout=10
+ Cload=1p
BGAIN NDRV 0 V=limit({Aol}*(V(INP)-V(INN)), V(VEE)+0.2, V(VCC)-0.2)
ROUT NDRV OUT {Rout}
COUT OUT 0 {Cload}
.ends CLAW_IDEAL_OPAMP
\end{verbatim}
}

Na versão compensada da rede com LT1810, o pós-processador
acrescenta, a cada estágio com realimentação, o capacitor de
compensação e o resistor de equilíbrio de polarização, este calculado
como o paralelo exato de todos os resistores do nó de soma do
estágio. O mesmo neurônio do trecho anterior torna-se:

{\small
\begin{verbatim}
XS4_0 bS4_0 s4_0 vcc veer a4_0 LT1810
RBALS4_0 bS4_0 0 8.89269k
RF4_0 a4_0 s4_0 100k
CFRF4_0 a4_0 s4_0 2p
   ... (ramos de peso inalterados)
XI4_0 bI4_0 si4_0 vcc vee na4_0 LT1810
RBALI4_0 bI4_0 0 50k
RIF4_0 na4_0 si4_0 100k
CFRIF4_0 na4_0 si4_0 2p
\end{verbatim}
}

\section{Código e reprodução}
\label{ap:codigo}

Os scripts residem em \texttt{experiments/analog-mlp/}:

\begin{itemize}
  \item \texttt{train\_analog\_mlp.py} -- obtém o MNIST, treina com as
        restrições do circuito e grava os pesos;
  \item \texttt{analog\_mlp.py} -- gera as três netlists principais,
        os esquemáticos e a figura de topologia;
  \item \texttt{bench\_741\_sat.cir} e
        \texttt{bench\_moderno\_sat.cir} -- calibram a saturação do
        LM741 e do LT1810; \texttt{bench\_vtc.cir} -- mede a curva de
        transferência final dos três amplificadores nos trilhos
        calibrados;
  \item \texttt{bench\_chain\_741.cir} e
        \texttt{bench\_chain\_moderno.cir} -- medem a latência do
        caminho crítico de cinco estágios;
  \item \texttt{analog\_mlp\_moderno\_fast.py} e
        \texttt{analog\_mlp\_moderno\_comp.py} -- geram as redes
        completas de janela reduzida com LT1810, sem e com as
        correções de compensação e equilíbrio;
  \item \texttt{analyze\_results.py} -- compara as simulações
        principais com o modelo, extrai latências e gera figuras e
        tabelas; \texttt{analyze\_moderno\_rede.py} -- analisa as
        redes de janela reduzida, a curva de transferência e o esquema
        de entrada, gerando as figuras e macros correspondentes;
  \item \texttt{bench\_host\_inference.py} -- mede o tempo de
        inferência em numpy no hospedeiro.
\end{itemize}

Sequência completa de reprodução:

\begin{verbatim}
./claw-spice code build experiments/analog-mlp/train_analog_mlp.py \
    --output-dir experiments/analog-mlp
./claw-spice code build experiments/analog-mlp/analog_mlp.py \
    --output-dir experiments/analog-mlp
./claw-spice code build experiments/analog-mlp/analog_mlp_moderno_fast.py \
    --output-dir experiments/analog-mlp
./claw-spice code build experiments/analog-mlp/analog_mlp_moderno_comp.py \
    --output-dir experiments/analog-mlp
./claw-spice sim run experiments/analog-mlp/analog_mlp.cir --timeout 2400
./claw-spice sim run experiments/analog-mlp/analog_mlp_741.cir --timeout 2400
./claw-spice sim run experiments/analog-mlp/bench_chain_741.cir --timeout 1800
./claw-spice sim run experiments/analog-mlp/bench_chain_moderno.cir --timeout 1800
./claw-spice sim run experiments/analog-mlp/bench_vtc.cir --timeout 600
./claw-spice sim run experiments/analog-mlp/analog_mlp_moderno_fast.cir \
    --timeout 2400
./claw-spice sim run experiments/analog-mlp/analog_mlp_moderno_comp.cir \
    --timeout 2400
./claw-spice code build experiments/analog-mlp/analyze_results.py \
    --output-dir experiments/analog-mlp
./claw-spice code build experiments/analog-mlp/analyze_moderno_rede.py \
    --output-dir experiments/analog-mlp
python3 experiments/analog-mlp/bench_host_inference.py
./claw-spice report analog-mlp
\end{verbatim}
